% ------------------------------------------------------------
% Appendix E — Spectral, Entropy, and Manifold Diagnostics
% ------------------------------------------------------------

\chapter{Spectral, Entropy, and Manifold Diagnostics}

This appendix presents the formal tools required to analyze the spectral, entropic, and geometric behavior of delta–sigma modulators when interpreted through the derived field-theoretic and manifold-based framework developed throughout the book. Although these tools are often invoked implicitly in engineering design, here they are constructed from first principles so that their relationship to the derived PDE, the entropy field, and the curvature of the underlying manifold becomes clear. Each diagnostic method provides a window into a different aspect of the modulator’s information-processing role, and collectively they form a complete set of invariants capturing the geometry of feedback-driven quantization.

\section{E.1 Power Spectral Density as a Derived Observable}

The power spectral density is usually introduced as a Fourier-domain representation of the second-order statistics of the error signal. However, in the geometric formulation the PSD is not merely a frequency-domain measure but rather a derived observable associated with the tangent complex of the modulator’s state manifold. The quantization error can be interpreted as a symbolic projection of the full PDE onto a discrete alphabet. This projection induces a nontrivial structure on the error manifold, and the PSD captures the manner in which the error energy distributes across the eigenmodes of the derived operator.

Let \(e[n]\) denote the quantization error. The PSD is defined as
\[
S_{e}(e^{j\omega}) = \lim_{N\to\infty} \frac{1}{N} \mathbb{E}\bigl|\sum_{n=0}^{N-1} e[n] e^{-j\omega n}\bigr|^{2}.
\]
In the geometric interpretation the Fourier modes correspond to eigenfunctions of the discrete Laplacian of the derived manifold. The noise-shaping property of the modulator corresponds to the annihilation of low-order derived modes. The zeros of the noise transfer function correspond to eigenmodes of the derived Laplacian that lie in the kernel of the symbolic projection. Thus each NTF zero is a geometric statement about the structure of the modulator’s tangent complex.

The PSD therefore becomes a diagnostic for curvature, symbolic collapse, and entropy routing. When the entropy potential is steep in low-frequency directions, the PSD decreases accordingly. When symbolic curvature is mild, the high-frequency tail of the PSD remains smooth. When derived homology classes persist, idle tones appear as persistent peaks in the spectrum. The PSD thus encodes global geometric properties in a form suitable for empirical measurement.

\section{E.2 Entropy Routing and the Distribution of Quantization Uncertainty}

Entropy routing is the process by which the modulator transports quantization uncertainty away from the signal band. In classical language this is the purpose of noise shaping. In the geometric formulation entropy routing is a well-defined vector field on the entropy manifold. Let the local entropy density be
\[
S_{e}(X) = -\log p(e|X),
\]
where \(e\) is the quantization error conditioned on the modulator state \(X\). The entropy current is defined as
\[
J_{e}(X) = -\nabla_{X} S_{e}(X) + w(X),
\]
where \(w(X)\) is the derived correction produced by the symbolic projection. The divergence of this current,
\[
\nabla\cdot J_{e}(X),
\]
measures the rate at which entropy is dissipated or accumulated in different regions of the manifold. Stable modulators exhibit negative divergence in the region where the input resides, meaning that entropy flows outward into the high-frequency band. Unstable modulators may exhibit positive divergence due to curvature collapse or excessive symbolic curvature, indicating that entropy is being trapped or accumulated in certain regions. This accumulation manifests as idle tones or overload phenomena.

By plotting or computing \(J_{e}(X)\) along simulated trajectories it is possible to detect subtle geometric defects in modulator design. A nearly vanishing divergence indicates efficient entropy transport. Strong local peaks indicate regions where the quantizer geometry and vector field misalign, producing undesirable harmonic content.

\section{E.3 Curvature Diagnostics on the State Manifold}

The manifold on which the modulator state evolves is endowed with a metric derived from the entropy potential. The curvature of this metric plays a fundamental role in stability, noise shaping, and spectral distribution. Let the Riemannian metric be
\[
g_{ij}(X) = \frac{\partial^{2} S}{\partial X^{i}\partial X^{j}}.
\]
The curvature tensor,
\[
R^{i}{}_{jkl},
\]
measures how geodesics deviate from one another. The scalar curvature,
\[
R = g^{jk} R^{i}{}_{jik},
\]
captures the average curvature at a point.

In the context of delta–sigma modulation negative curvature directions correspond to natural entropy-descent channels that drive error energy away from the signal band. Positive curvature directions indicate regions where the vector flow may slow or reverse, creating opportunities for quantization error to accumulate. The curvature invariants therefore provide a quantitative method for assessing whether a modulator’s loop filter and quantizer geometry produce the desired global flow.

Curvature diagnostics also detect the onset of nonlinear or chaotic regimes. When the curvature fluctuates rapidly over the manifold, the derived tangent complex experiences oscillatory growth, giving rise to quasi-periodic error modes. These modes are visible in the PSD as sidebands or harmonically related clusters.

\section{E.4 Divergence of the Vector Field and Stability Assessment}

The modulator’s stability is closely tied to the divergence of the vector field governing the flow. Let the flow be
\[
\frac{dX}{dt} = F(X).
\]
The divergence is
\[
\nabla\cdot F = \sum_{i} \frac{\partial F_{i}}{\partial X^{i}}.
\]
Negative divergence indicates volume contraction in state space, corresponding to stability through entropy descent. Positive divergence indicates local expansion, corresponding to instability or overload.

However, the symbolic projection imposed by the quantizer complicates this picture. The divergence must be evaluated on the derived manifold. The derived divergence is given by
\[
\mathrm{div}_{\mathrm{der}} F
= \mathrm{div} F - \mathrm{rank}\bigl(\mathrm{Ob}_{X}\bigr),
\]
where \(\mathrm{Ob}_{X}\) is the obstruction complex at the quantizer boundary. This subtraction reflects the fact that symbolic collapse removes certain tangent directions, effectively reducing the dimension of the flow. If the obstruction complex has high rank, the derived divergence may become positive even if the classical divergence is negative, signaling that too many stabilizing directions have been lost.

As a practical diagnostic, one computes the divergence at each timestep of a simulation and monitors its derived version when the state crosses quantizer thresholds. Patterns in the divergence indicate whether stability is globally preserved or lost in localized regions.

\section{E.5 Derived Harmonic Analysis and Idle-Tone Detection}

Idle tones are an important diagnostic in delta–sigma systems. They arise when the derived tangent complex possesses periodic or nearly periodic homology classes. In spectral terms idle tones appear as sharp peaks at specific frequencies. In geometric terms they arise when symbolic collapse and the vector field interact in such a way that the flow becomes trapped in a low-dimensional cycle of the derived manifold.

Derived harmonic analysis detects idle tones by studying the Laplacian of the derived manifold. The operator
\[
\Delta_{\mathrm{der}} = d d^{\ast} + d^{\ast} d,
\]
acts on the derived cochain complex. If \(\Delta_{\mathrm{der}}\) has low-lying eigenvalues corresponding to nontrivial derived homology, the associated eigenmodes produce periodic error oscillations. These oscillations manifest directly as idle tones.

By computing the spectrum of \(\Delta_{\mathrm{der}}\) numerically for a given modulator architecture, one can predict the presence of idle tones before hardware implementation. This diagnostic method also enables the design of dithering strategies that break the homology classes by smoothing the quantizer boundaries and thereby destroying the periodic derived structure.

\section{E.6 Summary of Diagnostic Integration}

The diagnostics developed in this appendix provide a unified view of delta–sigma behavior through spectral theory, entropy flow, curvature analysis, divergence measures, and derived harmonic invariants. Each diagnostic captures a different cross-section of the modulator’s geometry. Power spectral density reveals the global frequency distribution of error. Entropy routing exposes how uncertainty is transported across the manifold. Curvature identifies geometric regions favorable or unfavorable to dissipation. Divergence indicates stability or instability. Derived harmonic analysis detects persistent symbolic cycles.

Together these methods form a comprehensive diagnostic toolkit for analyzing, designing, and validating delta–sigma modulators in the geometric–derived–field-theoretic framework.

